% Consolidated theoretical manuscript; Paper A and Paper B remain preserved.
% Official user-provided Elsevier CAS 2.4 class: copy three files from
% paper/els-cas-templates/ into this directory before compilation.
\documentclass[a4paper,fleqn]{cas-dc}
\usepackage[authoryear,longnamesfirst]{natbib}
\usepackage{amsmath,amssymb,amsthm,mathtools}
\usepackage{graphicx,booktabs,tabularx,enumitem,url}
\usepackage{hyperref}
\hypersetup{colorlinks=true,urlcolor=blue,citecolor=blue,linkcolor=blue}
\newcommand{\R}{\mathbb R}
\newcommand{\tr}{\operatorname{tr}}
\newcommand{\diag}{\operatorname{diag}}
\newcommand{\RSS}{\operatorname{RSS}}
\newcommand{\edf}{\operatorname{edf}}
\newcommand{\CV}{\operatorname{CV}}
\newcommand{\GCV}{\operatorname{GCV}}
\newcommand{\op}{\operatorname{op}}
\newcommand{\1}{\mathbf 1}
\newtheorem{theorem}{Theorem}
\newtheorem{proposition}[theorem]{Proposition}
\newtheorem{corollary}[theorem]{Corollary}
\newtheorem{lemma}[theorem]{Lemma}
\newtheorem{remark}[theorem]{Remark}
\begin{document}
\let\WriteBookmarks\relax
\def\floatpagepagefraction{1}
\def\textpagefraction{.001}
\shorttitle{Algorithms for smoothness-indexed penalized least squares}
\shortauthors{Espino-Montelongo and Cort\'es-Toto}
\title[mode=title]{Spectral and structured algorithms for smoothness-indexed penalized least-squares trend estimation}
% Provisional author ordering, affiliations, declarations: coauthor review required.
\author[1]{Heriberto Espino-Montelongo}
\author[1]{Daniela Cort\'es-Toto}
\affiliation[1]{organization={Universidad de las Am\'ericas Puebla},
 city={San Andr\'es Cholula},state={Puebla},country={Mexico}}
\begin{abstract}
Finite-difference penalized least squares admits efficient structured solvers, yet choosing a smoothing level by cross-validation or an information criterion is a substantially different computational task from fitting one trend. We develop a unified computational treatment of this distinction using a normalized, trace-based smoothness coordinate on the closed unit interval. A singular-value representation of the arbitrary-order difference operator yields a strictly monotone smoothness-to-penalty map, its exact interpolation and constrained-polynomial limits, and a shift-free formula for a prescribed nonzero reference mean of differences. We identify the quantities required by conditional leave-one-out cross-validation, generalized cross-validation, information criteria and a Gaussian reduced-coordinate likelihood. We then derive and compare complete spectral and banded-Cholesky/selected-inverse pipelines, including their preprocessing, inverse smoothness mapping, criterion-specific scoring, memory requirements and reuse across candidate grids and observed series. A structured determinant identity avoids dense likelihood matrices. Exact arithmetic identities, conditional complexity bounds and a numerical-admissibility framework explain why the cost-minimizing algorithm depends on the entire workload: banded methods are natural for isolated fits and long series with small difference order, whereas a retained spectrum can amortize scalar scoring over many smoothness levels, subject to dense memory and conditioning constraints. An open-source Python interface exposes the formulation and the normalized coordinate. The contribution is an integrated theoretical and software-oriented framework; neither the classical difference penalty, trace index, spectral decomposition nor selected inversion is claimed as new.
\end{abstract}
\begin{keywords}
Penalized least squares \sep smoothness index \sep spectral decomposition
\sep Whittaker--Henderson smoothing \sep finite differences
\sep banded Cholesky \sep numerical stability \sep amortized complexity
\end{keywords}
\maketitle

\section{Introduction}\label{sec:introduction}
Finite-difference penalized least squares (PLS) is a classical approach to recovering a latent trend from noisy observations. Whittaker's quadratic-difference method \citep{whittaker1923}, the statistical formulation of \citet{guerrero2007}, and the broader Whittaker--Henderson tradition provide the foundation. A fixed nonzero mean of differences permits polynomial trends beyond the homogeneous special case. Guerrero's trace-based smoothness measure \citep{guerrero2008} and the comparison of classical optimality criteria by \citet{cortes2017} are likewise established starting points.

There are two different computational problems. First, at a given penalty $\lambda$, solve a structured positive-definite linear system to fit the trend. Second, when selecting a smoothing level, evaluate one or several criteria over many penalties. Cross-validation requires residuals and diagonal leverages, generalized cross-validation requires effective degrees of freedom, and a Gaussian likelihood needs quadratic forms and log determinants. Solving one banded system quickly does not establish the cost of evaluating an entire candidate grid. Conversely, paying for a dense spectral decomposition may be worthwhile when the result is reused.

The organizing variable in this paper is the user-facing normalized smoothness index $S\in[0,1]$. This reparameterization of the classical trace reduction is dimensionless and interpretable across orders: $S=0$ means interpolation; $S=1$ is the exactly constrained polynomial limit. A single SVD of the finite-difference operator supplies a monotone map between $S$ and the underlying penalty, fitted trends, and several candidate-scoring terms. The compact interval does not by itself guarantee globally optimal or faster numerical searches; its inversion has a cost, particularly near the endpoints.

The reference level $\mu$ is central to the formulation. The convenient assumption $\mu=0$ changes the limiting polynomial family and excludes prescribed nonzero $d$th-difference levels. While arbitrary fixed $\mu$ has already been treated by \citet{guerrero2007}, it is useful to have one numerically stable computational API handling both cases. In particular, explicitly forming a large degree-$d$ reference polynomial can lead to severe cancellation. Our shift-free singular-vector reconstruction avoids that intermediate.

The numerical and statistical foundations of efficient smoothing are extensive. \citet{weinert2007} exploits Whittaker--Henderson structure for efficient GCV, \citet{golub1979} develops classical spectral GCV methods, and \citet{biessy2026} supplies a modern Whittaker--Henderson statistical and computational treatment. We do not claim invention of diagonalization, the penalty, trace smoothness, cross-validation, Cholesky, or selected inversion. The contribution is a unified, reproducible \emph{workload-aware} treatment with explicit computational and numerical trade-offs, using the normalized smoothness coordinate throughout.

Specifically, the article develops (i) an arbitrary-order spectral formula supporting prescribed nonzero means without unstable polynomial shifts; (ii) the normalized index, its bijection and endpoints; (iii) the exact quantities required by conditional CV, GCV, information criteria and reduced Gaussian likelihood; (iv) time and storage bounds for complete spectral and banded pipelines across $M$ candidates, $R$ series and multiple criteria; and (v) a numerical-admissibility policy and open-source Python interface.

All principal theorems assume equally spaced observations, unweighted quadratic fidelity, integer $1\le d<n$, a fixed numeric $\mu$, and exact arithmetic. Data-derived reference means, time-series forecasting and fully measured hardware crossovers require additional statistical or empirical arguments, which are not silently incorporated into these claims.

\section{Finite-difference PLS and a spectral solution}\label{sec:model}
For equally spaced data $y\in\R^n$, let $1\le d<n$ and $r=n-d$. Define $D_d\in\R^{r\times n}$ by
\begin{equation}\label{eq:difference}
(D_dx)_i=\sum_{k=0}^d(-1)^{d-k}{d\choose k}x_{i+k}.
\end{equation}
At a fixed reference level $\mu\in\R$ of the $d$th differences, estimate
\begin{equation}\label{eq:pls}
\widehat\tau_{\lambda,\mu}=
\arg\min_{\tau\in\R^n}\{
\|y-\tau\|_2^2+\lambda\|D_d\tau-\mu\mathbf1_r\|_2^2\}.
\end{equation}
Set $Q_d=D_d^\top D_d$, $A_\lambda=I+\lambda Q_d$ and $H_\lambda=A_\lambda^{-1}$. The first-order equations give
\begin{equation}\label{eq:rhs}
A_\lambda\widehat\tau=y+\lambda\mu D_d^\top\mathbf1_r.
\end{equation}
$A_\lambda$ is symmetric positive definite for finite $\lambda\ge0$.

\begin{lemma}[Finite-difference nullspace]\label{lem:kernel}
$D_d$ has rank $r$ and nullspace dimension $d$.
Its nullspace consists of sampled polynomials of degree at most $d-1$; thus $Q_d$ has $d$ zero eigenvalues and $r$ strictly positive eigenvalues.
\end{lemma}
\begin{proof}
Repeated differences of a polynomial of degree below $d$ vanish. Conversely, if $D_dx=0$, induction on the difference order implies that $x_t$ is a polynomial in $t-1$ of degree at most $d-1$. The $d$ sequences $\binom{t-1}{k}$, $k=0,\ldots,d-1$, are independent. Rank--nullity and $x^\top Q_dx=\|D_dx\|^2$ complete the proof.
\end{proof}

The particular polynomial $p_{\mu,t}=\mu\binom{t-1}{d}$ satisfies $D_dp_\mu=\mu\mathbf1_r$. Hence
\begin{equation}\label{eq:affine}
\widehat\tau_{\lambda,\mu}=p_\mu+H_\lambda(y-p_\mu).
\end{equation}
For large $n,d$, explicitly forming and subtracting $p_\mu$ can cause severe cancellation. We therefore use the reduced SVD $D_d=U\Sigma V^\top$, where $\Sigma=\diag(\sigma_1,\ldots,\sigma_r)$, $\sigma_j>0$, $U$ is $r$-square orthogonal and $V$ has $r$ orthonormal columns. Define $\delta_j=\sigma_j^2$ and $w_j(\lambda)=\lambda\delta_j/(1+\lambda\delta_j)$.

\begin{proposition}[Shift-free nonzero-reference reconstruction]\label{prop:shiftfree}
For every finite $\lambda\ge0$ and specified real $\mu$,
\begin{equation}\label{eq:shiftfree}
\begin{aligned}
\widehat\tau_{\lambda,\mu}
&=y-V\bigl[w\odot(V^\top y)\\
&\qquad-\mu(w\oslash\sigma)\odot(U^\top\mathbf1_r)\bigr].
\end{aligned}
\end{equation}
Componentwise multiplication and division are denoted by $\odot$ and $\oslash$. This identity avoids building $p_\mu$ or inverting a dense matrix explicitly.
\end{proposition}
\begin{proof}
Since $H_\lambda=I-V\diag(w)V^\top$ and $D_dp_\mu=\mu\mathbf1_r$, we have $V^\top p_\mu=\mu\Sigma^{-1}U^\top\mathbf1_r$. Substitute into \eqref{eq:affine}.
\end{proof}

\begin{corollary}[Exact limiting fits]\label{cor:endfits}
At $\lambda=0$, $\widehat\tau=y$. As $\lambda\to\infty$,
\begin{equation}\label{eq:inftyfit}
\begin{aligned}
\widehat\tau_{\infty,\mu}
&=y-V\bigl[V^\top y-\mu\Sigma^{-1}U^\top\mathbf1_r\bigr]\\
&=\arg\min_{\tau:\,D_d\tau=\mu\mathbf1_r}\|y-\tau\|_2^2.
\end{aligned}
\end{equation}
The limiting fit has degree at most $d-1$ for $\mu=0$, and may have degree $d$ for fixed $\mu\ne0$.
\end{corollary}
\begin{proof}
The weights tend to zero and one at the respective limits. The second expression satisfies the affine constraint and removes only the component in the row space of $D_d$, so it is the orthogonal constrained projection.
\end{proof}

Fixed nonzero $\mu$ is already part of Guerrero's model \citep{guerrero2007}; it is not an invention of this paper. A data-derived plug-in such as $\widehat\mu=\mathbf1_r^\top D_dy/r$ changes the observation-to-fit derivative and requires different treatment of exact leave-one-out refitting. Unless stated otherwise our results are conditional on a fixed numeric $\mu$.

\section{A normalized and invertible smoothness coordinate}\label{sec:smoothness}
For a fixed numeric $\mu$, define the conditional effective degrees of freedom $\nu_d(\lambda)=\tr H_\lambda$. The $d$ polynomial-nullspace eigenvalues of $H_\lambda$ equal one, and the other $r=n-d$ eigenvalues equal $(1+\lambda\delta_j)^{-1}$. Therefore
\begin{equation}\label{eq:edf}
\nu_d(\lambda)=d+\sum_{j=1}^r\frac{1}{1+\lambda\delta_j}.
\end{equation}
The classical trace reduction divided by $n$ is
\begin{equation}\label{eq:classicalS}
S^{\rm trace}_d(\lambda)=\frac{n-\nu_d(\lambda)}{n},
\quad 0\le S^{\rm trace}_d\le1-\frac dn.
\end{equation}
The normalized coordinate used here divides by the number of \emph{penalized} directions:
\begin{equation}\label{eq:S}
\boxed{\;
S_d(\lambda)=\frac{n-\nu_d(\lambda)}{n-d}
=\frac{1}{r}\sum_{j=1}^r\frac{\lambda\delta_j}{1+\lambda\delta_j}.
\;}
\end{equation}
Thus $S_d=(n/r)S_d^{\rm trace}$, an interpretable rescaling of a published index \citep{guerrero2008,cortes2017}, not a different family of fitted trends.

\begin{theorem}[Smoothness bijection]\label{thm:bijection}
For every integer $1\le d<n$, $S_d$ is continuous and strictly increasing on the nonnegative finite real line, with $S_d(0)=0$ and $\lim_{\lambda\to\infty}S_d(\lambda)=1$. It extends to a bijection between $[0,\infty]$ and $[0,1]$. Each $s\in(0,1)$ has a unique inverse $\lambda=\Lambda_{n,d}(s)$ independent of $y$ and any fixed numeric $\mu$.
\end{theorem}
\begin{proof}
All $\delta_j$ are strictly positive by Lemma~\ref{lem:kernel}, so
\begin{equation}\label{eq:Sprime}
S_d'(\lambda)=\frac{1}{r}\sum_{j=1}^r
\frac{\delta_j}{(1+\lambda\delta_j)^2}>0.
\end{equation}
Continuity and the stated endpoint values follow term by term. The inverse exists by the intermediate value theorem and strict monotonicity. Formula \eqref{eq:S} contains only $(n,d,\lambda)$ and the spectrum of $D_d$, so it is independent of the observed values and fixed reference mean.
\end{proof}

The same requested $s$ has a common interpretation across different orders, but generally corresponds to \emph{different} penalties $\Lambda_{n,d}(s)$. Direct optimization in raw $\lambda$ or $\log\lambda$ is possible. The point of $s$ is a bounded, physically interpretable search parameter, not a proof of optimization supremacy.

\begin{proposition}[Change of selection coordinate]\label{prop:coordinate}
Let $C(\lambda)$ be any fixed-$\mu$ criterion on an admissible finite interval and set $\widetilde C(s)=C(\Lambda_{n,d}(s))$. The two objectives attain the same global minimizers after the monotone change of coordinates, provided the domains and endpoint conventions match. For $s\in(0,1)$,
\begin{equation}\label{eq:chain}
\frac{d\widetilde C}{ds}
=\frac{dC/d\log\lambda}{dS_d/d\log\lambda},
\quad
\frac{dS_d}{d\log\lambda}
=\frac1r\sum_{j=1}^r\frac{\lambda\delta_j}
{(1+\lambda\delta_j)^2}.
\end{equation}
\end{proposition}
\begin{proof}
The first statement follows from the bijection. The second is the chain rule and $d\lambda/d\log\lambda=\lambda$.
\end{proof}
The denominator in \eqref{eq:chain} tends to zero at both endpoints; hence boundedness alone does not resolve numerical conditioning or spurious local minima of the criterion. Global minimizer equivalence does not imply that bounded and logarithmic numerical optimizers follow identical paths.

\paragraph{Exact endpoint conventions.}
The point $s=0$ means $\lambda=0$ and the interpolation fit $y$; the point $s=1$ means the constrained fit \eqref{eq:inftyfit}, not a floating-point stand-in for an extremely large finite penalty. Some criteria are undefined at these endpoints even though the fitted trends exist. A robust implementation must treat such limits analytically rather than blindly calling a root finder at $s=0$ or $s=1$.

\paragraph{Caching the inverse map.}
For fixed $n,d$, the full set $\{\Lambda_{n,d}(s_j)\}_{j=1}^M$ can be mapped once and reused across multiple fixed-$\mu$ values, independent series, and criteria. In contrast, an independently constructed uniform grid in $\log\lambda$ generally selects different candidate trends and is not an identical-workload control.

\section{Complete statistical criteria}\label{sec:criteria}
The cost of \emph{selecting} a smoothing level is not the cost of one fit. For fixed numeric $\mu$ and finite $\lambda>0$, write $e_\lambda=y-\widehat\tau_{\lambda,\mu}$, $\RSS_\lambda=e_\lambda^\top e_\lambda$, $h_i=(H_\lambda)_{ii}$, and $\nu=\tr H_\lambda$. The exact classical leave-one-out and generalized cross-validation scores are
\begin{equation}\label{eq:cvgcv}
\CV(\lambda)=\frac1n\sum_{i=1}^n
\left(\frac{e_{\lambda,i}}{1-h_i}\right)^2,\qquad
\GCV(\lambda)=\frac{\RSS_\lambda/n}{(1-\nu/n)^2}.
\end{equation}
These are retrospective smooth-trend criteria, not chronological forecasting-validation scores. The fixed-$\mu$ leverage diagonal is available spectrally as
\begin{equation}\label{eq:leverage}
h_i=1-\sum_{j=1}^r V_{ij}^2 w_j(\lambda).
\end{equation}
The unpenalized polynomial projection is implicitly included by the leading one in \eqref{eq:leverage}. Omitting it would give incorrect CV denominators.

For comparability with the conventions in \citet{cortes2017}, we also consider the residual-based expressions
\begin{equation}\label{eq:ic}
\begin{aligned}
\mathrm{AICc}_{\RSS}(\lambda)
 &=\log(\RSS_\lambda/n)+\frac{2\nu+1}{n-\nu-2},\\
\mathrm{BIC}_{\RSS}(\lambda)
 &=\log(\RSS_\lambda/n)+\frac{\nu\log n}{n}.
\end{aligned}
\end{equation}
The AICc value is defined only when its denominator is admissible, and the names here refer to precisely the stated residual-based conventions, not to every information criterion with the same name in other statistical models.

\begin{proposition}[Residual-BIC interpolation pathology]
\label{prop:bic}
For fixed $\mu$ suppose $Q_d(y-p_\mu)\ne0$. Then the expression $\mathrm{BIC}_{\RSS}(\lambda)$ in \eqref{eq:ic} tends to $-\infty$ as $\lambda\downarrow0$, and hence as $s=S_d(\lambda)\downarrow0$.
\end{proposition}
\begin{proof}
Since $H_\lambda=I-\lambda Q_d+O(\lambda^2)$,
$e_\lambda=(I-H_\lambda)(y-p_\mu)=
\lambda Q_d(y-p_\mu)+O(\lambda^2)$.
Thus $\RSS_\lambda=\lambda^2\|Q_d(y-p_\mu)\|^2+O(\lambda^3)$, while $\nu\to n$. The logarithmic residual term equals $2\log\lambda+O(1)$ and dominates the finite degrees-of-freedom penalty. The smoothness limit follows from Theorem~\ref{thm:bijection}.
\end{proof}
A strictly positive imposed grid minimum may be a disclosed sensitivity restriction, but is not a repaired global optimum for this expression. Likewise, CV/GCV near exact interpolation require careful limiting conventions rather than evaluating an indeterminate ratio at $s=0$.

\subsection{Gaussian reduced-coordinate likelihood}
To distinguish statistical selection from penalized training residuals, suppose in addition that
\begin{equation}\label{eq:gaussian}
y=\tau+\varepsilon,\quad
\varepsilon\sim N(0,\sigma_\varepsilon^2I_n),\quad
D_d\tau-\mu\mathbf1_r\sim N(0,\sigma_u^2I_r),
\quad \lambda=\sigma_\varepsilon^2/\sigma_u^2,
\end{equation}
with independent components and a fixed numeric $\mu$. The $r$ penalized SVD coordinates $z=V^\top(y-p_\mu)$ have marginal variances $\sigma_\varepsilon^2c_j(\lambda)$, where $c_j=1+1/(\lambda\delta_j)$. Profiling $\sigma_\varepsilon^2$ gives the criterion
\begin{equation}\label{eq:ml}
\mathcal L(\lambda)=
r\log\left(\frac{1}{r}\sum_{j=1}^r
\frac{z_j^2}{c_j(\lambda)}\right)
+\sum_{j=1}^r\log c_j(\lambda),
\end{equation}
up to a constant independent of $\lambda$. It is a reduced-coordinate variance-component likelihood, not the result of minimizing the PLS training objective with respect to $\lambda$.

\begin{proposition}[Banded likelihood sufficiency]\label{prop:det}
Let $\lambda>0$, $B_d=D_dD_d^\top$, and let $\widehat\tau$ solve \eqref{eq:rhs}. Set
$J_\lambda=\|y-\widehat\tau\|^2+
\lambda\|D_d\widehat\tau-\mu\mathbf1_r\|^2$.
Then
\begin{equation}\label{eq:quadraticidentity}
J_\lambda=\sum_{j=1}^r\frac{z_j^2}{c_j(\lambda)},
\end{equation}
and
\begin{equation}\label{eq:detidentity}
\sum_{j=1}^r\log c_j(\lambda)
=\log\det A_\lambda-r\log\lambda-\log\det B_d.
\end{equation}
Consequently a banded factorization of $A_\lambda$ and a precomputed banded factorization of $B_d$ suffice to evaluate \eqref{eq:ml} without a dense inverse.
\end{proposition}
\begin{proof}
In the positive spectral directions the minimal value of the quadratic PLS objective is $\sum_j w_j z_j^2$. Since $w_j=c_j^{-1}$ this proves \eqref{eq:quadraticidentity}; the nullspace part of the objective vanishes at the optimum. The positive eigenvalues of $Q_d$ and $B_d$ coincide. Thus
$\log\det A_\lambda=\sum_j\log(1+\lambda\delta_j)$ and $\log\det B_d=\sum_j\log\delta_j$. Subtracting $r\log\lambda$ proves \eqref{eq:detidentity}. Both determinants admit banded Cholesky log-diagonal evaluations.
\end{proof}

The same likelihood formula must not be called joint unknown-$\mu$ maximum likelihood after merely substituting $\widehat\mu=\operatorname{mean}(D_dy)$. Such a plug-in policy changes the statistical estimation problem. The direct PLS fit, likelihood fit, and genuinely foldwise-refitted leave-one-out selector are distinct objects.

\section{Numerical algorithm families}\label{sec:algorithms}
The principal distinction is between a fitted trend at a \emph{known} penalty and a \emph{complete criterion evaluation} at that penalty. We include every required residual, trace, leverage, determinant and smoothness-mapping computation in the latter.

\subsection{Spectral preparation and candidate scoring}
A full retained reduced SVD of $D_d$ requires $O(nr^2)$ arithmetic and $O(nr)$ memory for $r=n-d$ (near-square: cubic time and quadratic storage) under a standard dense bidiagonalization-based SVD computation \citep{golubvanloan2013}. The stated count is for a general dense rectangular factorization, not a lower bound on all solvers exploiting $D_d Preparation produces both the positive singular spectrum and the right singular vectors. For one observed series, projections $V^\top y$ and $U^\top\mathbf1_r$ permit scalar evaluations of $S$, $\nu$, $\RSS$, GCV, residual AICc/BIC and the Gaussian criterion using $O(r)$ arithmetic per candidate after an $O(nr)$ projection. A full curve reconstructs via \eqref{eq:shiftfree} in $O(nr)$. Ordinary CV additionally computes the leverage vector \eqref{eq:leverage} and reconstructed residuals, each up to $O(nr)$ per new penalty. Leverages can be shared across series having the same $(n,d,\lambda)$.

Dense eigenanalysis of $Q_d$ supplies comparable spectral formulas but explicitly forming $D_d^\top D_d$ squares the nonzero singular-value condition number \citep{golubvanloan2013}. A full SVD of $D_d$ is the reference accuracy-oriented spectral construction in our Python implementation; it is not claimed to minimize runtime.

\subsection{Banded Cholesky and selected inversion}
$A_\lambda$ is SPD and has half-bandwidth $d$. Factor $A_\lambda=L_\lambda L_\lambda^\top$ with lower triangular banded $L_\lambda$, in $O(nd^2)$ arithmetic and $O(nd)$ storage; solve \eqref{eq:rhs} using two triangular solves, $O(nd)$ per right-hand side. These are classical banded-direct bounds \citep{golubvanloan2013,davis2006}; efficient PLS-specific structured and GCV algorithms precede this work \citep{weinert2007}. The direct right-hand side avoids forming the high-degree polynomial $p_\mu$. For just one fit, this is asymptotically preferable to preparing a dense SVD when $d$ is fixed and $n$ grows.

However, exact CV requires diagonal entries of $Z=A_\lambda^{-1}$. Let $L=L_\lambda$ and denote the selected band of $Z$ by $Z_{ij}$ for $|i-j|\le d$. For $j>i$, backward selected-inverse recurrences are
\begin{equation}\label{eq:seloff}
Z_{ij}=-\frac1{L_{ii}}\sum_{k=i+1}^{\min(n,i+d)}
 L_{ki}Z_{kj},
\qquad i<j\le\min(n,i+d),
\end{equation}
and
\begin{equation}\label{eq:seldiag}
Z_{ii}=\frac{1}{L_{ii}^2}
-\frac{1}{L_{ii}}\sum_{k=i+1}^{\min(n,i+d)}
L_{ki}Z_{ki}.
\end{equation}
By symmetry $Z_{kj}=Z_{jk}$. The off-diagonal recursion follows by expanding $(L^\top Z)_{ij}=0$ for $i<j$, and the diagonal relation from $(L^\top Z)_{ii}=1/L_{ii}$. They are classical selected-inverse identities \citep{erisman1975}, not an original inverse algorithm. The $O(nd^2)$ bound below follows by counting $O(d)$ updates for each of $O(nd)$ selected entries; for fixed $d$, it reduces to $O(n)$, consistent with the specialized Whittaker--Henderson algorithms in \citet{weinert2007}.

The selected entries are computable in reverse order in $O(nd^2)$ work and $O(nd)$ storage. Their diagonal supplies $h_i=Z_{ii}$ and $\nu=\sum_i h_i$. Consequently exact conditional CV, GCV and information criteria do \emph{not} require solving for all $n$ columns of the inverse or forming the dense $H_\lambda$. The determinant identity in Proposition~\ref{prop:det} completes the same banded pipeline for the Gaussian likelihood.

\subsection{Other exact and iterative approaches}
General banded LU can solve the finite-penalty system at structured cost but is not SPD Cholesky. Sparse direct solvers can exploit the same graph structure; the particular package, pivoting, fill-in, ordering and factorization type must be named correctly \citep{davis2006}. A hybrid in which sparse LU or CHOLMOD supplies a trend but a \emph{separate} banded factorization supplies leverage or likelihood diagnostics must charge both steps.

Matrix-free preconditioned conjugate gradients (PCG) applies $A_\lambda$ through local differences at $O(nd)$ work per iteration, producing a nominal $O(qnd)$ fit for $q$ iterations, plus preconditioner setup and application \citep{saad2003}. The $O(nd)$ matrix-vector application is counted directly from the two order-$d$ difference passes; the iteration count $q$ depends on convergence and the preconditioned spectrum. FFT or sine-transform preconditioners can introduce $O(qn\log n)$ work. Approximate boundary transforms should not be described as an exact finite-interval eigendecomposition. Neither a universally bounded iteration count nor exact leverage access follows from PCG alone; convergence and added diagnostics are separate requirements.

\begin{table*}[t]
\caption{Representative leading computational costs at one known finite $\lambda$. Standard factorization counts follow \citet{golubvanloan2013,davis2006}; selected-inverse and PLS-specialized counts follow \citet{erisman1975,weinert2007}; iterative costs follow \citet{saad2003}. For full spectral methods the decomposition is reusable. The full conditional CV score includes leverage and residual computation. These are arithmetic bounds, not measured seconds.}
\label{tab:algorithms}\centering\small
\setlength{\tabcolsep}{4pt}
\begin{tabular}{@{}p{.25\textwidth}p{.22\textwidth}p{.23\textwidth}p{.16\textwidth}@{}}
\toprule
Pipeline & Preparation / factorization & Candidate fit or complete score & Working memory\\
\midrule
Dense inverse/solve & $O(n^3)$ per new $\lambda$ & $O(n^3)$ complete & $O(n^2)$\\
Full SVD of $D_d$ & $O(nr^2)$ once & $O(r)$ scalar score; $O(nr)$ full fit/CV & $O(nr)$\\
Dense Gram eigenvectors & $O(n^3)$ once & $O(n)$ scalar score; $O(n^2)$ full fit/CV & $O(n^2)$\\
Banded Cholesky & $O(nd^2)$ per $\lambda$ & $O(nd)$ one fit; diagnostics extra & $O(nd)$\\
Banded Cholesky with selected inverse & $O(nd^2)$ per $\lambda$ & $O(nd)$ solve; $O(nd^2)$ complete CV & $O(nd)$\\
PCG, $q$ iterations & preconditioner-dependent & $O(qnd)$ fit; exact criteria extra & method-dependent\\
\bottomrule
\end{tabular}
\end{table*}
s additional band structure. Preparation produces both the positive singular spectrum and the right singular vectors. For one observed series, projections $V^\top y$ and $U^\top\mathbf1_r$ permit scalar evaluations of $S$, $\nu$, $\RSS$, GCV, residual AICc/BIC and the Gaussian criterion using $O(r)$ arithmetic per candidate after an $O(nr)$ projection. A full curve reconstructs via \eqref{eq:shiftfree} in $O(nr)$. Ordinary CV additionally computes the leverage vector \eqref{eq:leverage} and reconstructed residuals, each up to $O(nr)$ per new penalty. Leverages can be shared across series having the same $(n,d,\lambda)$.

Dense eigenanalysis of $Q_d$ supplies comparable spectral formulas but explicitly forming $D_d^\top D_d$ squares the nonzero singular-value condition number. A full SVD of $D_d$ is the reference accuracy-oriented spectral construction in our Python implementation; it is not claimed to minimize runtime.

\subsection{Banded Cholesky and selected inversion}
$A_\lambda$ is SPD and has half-bandwidth $d$. Factor $A_\lambda=L_\lambda L_\lambda^\top$ with lower triangular banded $L_\lambda$, in $O(nd^2)$ arithmetic and $O(nd)$ storage; solve \eqref{eq:rhs} using two triangular solves, $O(nd)$ per right-hand side. The direct right-hand side avoids forming the high-degree polynomial $p_\mu$. For just one fit, this is asymptotically preferable to preparing a dense SVD when $d$ is fixed and $n$ grows.

However, exact CV requires diagonal entries of $Z=A_\lambda^{-1}$. Let $L=L_\lambda$ and denote the selected band of $Z$ by $Z_{ij}$ for $|i-j|\le d$. For $j>i$, backward selected-inverse recurrences are
\begin{equation}\label{eq:seloff}
Z_{ij}=-\frac1{L_{ii}}\sum_{k=i+1}^{\min(n,i+d)}
 L_{ki}Z_{kj},
\qquad i<j\le\min(n,i+d),
\end{equation}
and
\begin{equation}\label{eq:seldiag}
Z_{ii}=\frac{1}{L_{ii}^2}
-\frac{1}{L_{ii}}\sum_{k=i+1}^{\min(n,i+d)}
L_{ki}Z_{ki}.
\end{equation}
By symmetry $Z_{kj}=Z_{jk}$. The off-diagonal recursion follows by expanding $(L^\top Z)_{ij}=0$ for $i<j$, and the diagonal relation from $(L^\top Z)_{ii}=1/L_{ii}$. They are classical selected-inverse identities \citep{erisman1975}, not an original inverse algorithm.

The selected entries are computable in reverse order in $O(nd^2)$ work and $O(nd)$ storage. Their diagonal supplies $h_i=Z_{ii}$ and $\nu=\sum_i h_i$. Consequently exact conditional CV, GCV and information criteria do \emph{not} require solving for all $n$ columns of the inverse or forming the dense $H_\lambda$. The determinant identity in Proposition~\ref{prop:det} completes the same banded pipeline for the Gaussian likelihood.

\subsection{Other exact and iterative approaches}
General banded LU can solve the finite-penalty system at structured cost but is not SPD Cholesky. Sparse direct solvers can exploit the same graph structure; the particular package and factorization type must be named correctly. A hybrid in which sparse LU or CHOLMOD supplies a trend but a \emph{separate} banded factorization supplies leverage or likelihood diagnostics must charge both steps.

Matrix-free preconditioned conjugate gradients (PCG) applies $A_\lambda$ through local differences at $O(nd)$ work per iteration, producing a nominal $O(qnd)$ fit for $q$ iterations, plus preconditioner setup and application. FFT or sine-transform preconditioners can introduce $O(qn\log n)$ work. Approximate boundary transforms should not be described as an exact finite-interval eigendecomposition. Neither a universally bounded iteration count nor exact leverage access follows from PCG alone; convergence and added diagnostics are separate requirements.

\begin{table*}[t]
\caption{Representative leading computational costs at one known finite $\lambda$. For full spectral methods the decomposition is reusable. The full conditional CV score includes leverage and residual computation. These are arithmetic bounds, not measured seconds.}
\label{tab:algorithms}\centering\small
\setlength{\tabcolsep}{4pt}
\begin{tabular}{@{}p{.25\textwidth}p{.22\textwidth}p{.23\textwidth}p{.16\textwidth}@{}}
\toprule
Pipeline & Preparation / factorization & Candidate fit or complete score & Working memory\\
\midrule
Dense inverse/solve & $O(n^3)$ per new $\lambda$ & $O(n^3)$ complete & $O(n^2)$\\
Full SVD of $D_d$ & $O(nr^2)$ once & $O(r)$ scalar score; $O(nr)$ full fit/CV & $O(nr)$\\
Dense Gram eigenvectors & $O(n^3)$ once & $O(n)$ scalar score; $O(n^2)$ full fit/CV & $O(n^2)$\\
Banded Cholesky & $O(nd^2)$ per $\lambda$ & $O(nd)$ one fit; diagnostics extra & $O(nd)$\\
Banded Cholesky with selected inverse & $O(nd^2)$ per $\lambda$ & $O(nd)$ solve; $O(nd^2)$ complete CV & $O(nd)$\\
PCG, $q$ iterations & preconditioner-dependent & $O(qnd)$ fit; exact criteria extra & method-dependent\\
\bottomrule
\end{tabular}
\end{table*}

\section{Amortized time and memory complexity}\label{sec:complexity}
Fix a length $n$, difference order $d$, an interior smoothness grid $\mathcal S=\{s_1,\ldots,s_M\}\subset(0,1)$, $R$ independent observed series of the same length, and a criterion set $\mathcal C$ with $c=|\mathcal C|$. Let $I$ be an upper bound on the scalar function evaluations used to invert each $s_j$ to the requested tolerance. This is a \emph{workload}, not merely one linear-system solve. Its total execution time decomposes conceptually as
\begin{equation}\label{eq:costpartition}
T_{\mathcal A}=T_{\rm prep}+T_{\rm map}
 +T_{\rm factor}+T_{\rm series}+T_{\rm scores}+T_{\rm overhead},
\end{equation}
where $\mathcal A$ is a particular complete algorithm and its configuration. The component costs use classical dense and banded linear algebra \citep{golubvanloan2013,davis2006}, classical selected-inverse identities \citep{erisman1975}, and pre-existing structured Whittaker--Henderson/GCV algorithms \citep{weinert2007}. We derive the following \emph{combined workload-specific upper bounds}; these formulas are not directly quoted from those sources, nor do they imply asymptotic optimality among all possible algorithms.

\begin{theorem}[Work for a retained spectral pipeline]\label{thm:svdwork}
Assume the full reduced SVD is retained; each interior smoothness inversion requires at most $I$ evaluations of the $r$-term spectral sum; each observed series is projected once; and projected scalar scores use $O(r)$ work. Put $\chi_{\rm CV}=1$ when ordinary conditional CV is requested and zero otherwise. An upper bound on arithmetic work is
\begin{equation}\label{eq:svdwork}
W_{\rm SVD}
=O(nr^2+MIr+Rnr+MRcr+\chi_{\rm CV}MRnr).
\end{equation}
Working storage is $O(nr+M+Rn)$ if all $R$ input series are kept and only best-score summaries are retained.
\end{theorem}
\begin{proof}
For a standard dense reduced SVD of an $r\times n$ matrix, $O(nr^2)$ arithmetic and $O(nr)$ retained storage follow from classical dense factorization bounds \citep{golubvanloan2013}. Each of the $MI$ smoothness evaluations sums $r$ positive-mode terms. One projection per observed series costs $O(nr)$. At each of $MR$ candidate-series pairs, each scalar score takes $O(r)$ work after projection. When CV is requested, the leverage diagonal and fitted residual vector can each be formed using $O(nr)$ matrix-vector work; multiplying by $MR$ bounds this additional cost. Caching leverages across series can improve constants but is not needed for the claimed bound. Candidate values and $R$ input series add $O(M+Rn)$ memory.
\end{proof}

\begin{theorem}[Work for a streamed banded pipeline]\label{thm:bandwork}
Assume exact SPD banded Cholesky and selected inversion are used both for smoothness-to-penalty inversion and for complete score diagnostics, and the resulting penalty grid is shared across criteria and series. Then
\begin{equation}\label{eq:bandwork}
W_{\rm band}
=O(nd^2+MInd^2+Mnd^2+MRnd+MRcn).
\end{equation}
A candidate-streamed implementation retaining all input series uses $O(nd+M+Rn)$ working storage. Keeping every candidate factorization instead requires $O(Mnd)$ extra factor storage.
\end{theorem}
\begin{proof}
Assembling the banded penalty and precomputing the Gaussian constant determinant cost at most $O(nd^2)$. A banded factorization and selected inverse cost $O(nd^2)$ per trial value, using standard banded-direct and selected-inverse algorithms \citep{golubvanloan2013,erisman1975,weinert2007}. Up to $MI$ trials yield the second term; factorizing the $M$ accepted penalties yields the third. For each candidate and series, banded triangular solves and difference passes cost $O(nd)$, and forming each criterion's scalar contributions costs at most $O(n)$. One factor and selected band are sufficient while processing all $R$ series at the current candidate. The storage claims follow by retaining or discarding these structures.
\end{proof}

\paragraph{Bounds versus measured crossover.}
The displayed $O(\cdot)$ expressions are upper bounds for specific computational families, not universal lower bounds on all possible implementations. They omit machine constants, cache contention and implementation-specific allocation; an optimistic operation count is not observed runtime. A solver can be admissible only after its accuracy and memory requirements are satisfied.

\begin{proposition}[Conditional candidate-count crossover]\label{prop:crossover}
Suppose two locally calibrated complete-pipeline times for fixed $(n,d,R,\mathcal C)$ admit
\[
T_s(M)=a_s+b_sM,\qquad T_b(M)=a_b+b_bM.
\]
If $a_s>a_b$ and $b_b>b_s$, equality occurs at the positive threshold
\begin{equation}\label{eq:Mstar}
M_\star=\frac{a_s-a_b}{b_b-b_s}.
\end{equation}
Banded is faster below this threshold and spectral above it under the stated local affine model. Without both sign conditions the claimed direction of crossover does not follow.
\end{proposition}
\begin{proof}
Subtract the affine functions to obtain $(a_s-a_b)-M(b_b-b_s)$ and inspect its sign.
\end{proof}

This proposition does not infer a numerical $M_\star$ from the preceding big-$O$ statements. In particular the $O(nr)$ spectral CV term can reverse a scalar-GCV comparison; repeated series lower effective setup cost per series; and $S\mapsto\lambda$ inversion may dominate an otherwise inexpensive one-fit banded approach.

\subsection{Special structure and hybrid mapping}
For first differences the penalty is the finite path-graph Laplacian. Its positive eigenvalues have the known analytic expression
\begin{equation}\label{eq:d1spectrum}
\delta_j=2-2\cos\left(\frac{\pi j}{n}\right),
\qquad j=1,\ldots,n-1.
\end{equation}
They can be generated using $O(n)$ arithmetic and storage without full SVD. The requested $M$ smoothness levels can consequently be mapped in $O(MIn)$ arithmetic using \eqref{eq:S}. This established discrete-cosine identity should not be generalized naively to arbitrary higher-order finite-interval differences: boundary effects change the spectrum.

More generally the mapping engine need not equal the fitting engine. If the spectrum (or a suitably certified approximation) is available, one may map $\mathcal S$ spectrally and then fit the resulting penalties with banded Cholesky. Such \emph{hybrid mapping} can retain the linear-memory final solve, but its own spectral preparation, error and storage must be accounted for. A dense full SVD built only to obtain one banded fitted trend may be a poor trade-off.

\paragraph{Memory frontier.}
A single float64 dense $n\times n$ array requires $8n^2$ bytes before copying or factoring, while fixed-small-$d$ banded storage grows linearly. At $n=131072$, one dense array alone occupies $128$ GiB. Therefore an inadmissible dense method must be rejected by a memory guard rather than reported as a completed comparison. When multiple series or smoothing candidates are processed, intermediate dense matrices should not be retained without an explicit justification.

\section{Numerical conditioning, accuracy and failure mechanisms}
\label{sec:stability}
An exact SPD factorization exists for all finite $\lambda$, but finite-precision arithmetic can still produce unusable fitted trends or selection scores. In particular, the difference order, penalty magnitude, extreme smoothness levels and criterion diagnostics affect numerical validity.

\begin{proposition}[Conditioning of the primal system]\label{prop:conditioning}
For $1\le d<n$ and finite $\lambda\ge0$,
\begin{equation}\label{eq:kappa}
\kappa_2(A_\lambda)=1+\lambda\delta_{\max}
\le 1+\lambda4^d,
\end{equation}
where $\delta_{\max}$ is the largest positive eigenvalue of $D_d^\top D_d$.
\end{proposition}
\begin{proof}
By Lemma~\ref{lem:kernel}, $A_\lambda$ has $d$ eigenvalues equal to one and remaining eigenvalues $1+\lambda\delta_j$. The ratio of its largest and smallest eigenvalues is $1+\lambda\delta_{\max}$. Each forward-difference stencil has absolute coefficient sum $2^d$, so $\|D_d\|_\infty\le2^d$ and $\|D_d\|_1\le2^d$. Consequently $\|D_d\|_2\le\sqrt{\|D_d\|_1\|D_d\|_\infty}\le2^d$, yielding $\delta_{\max}=\|D_d\|_2^2\le4^d$.
\end{proof}

The condition bound becomes demanding near $S=1$, for which $\lambda$ diverges. Forming the normal matrix explicitly also squares the nonzero singular-value condition number relative to working directly with $D_d$. These facts motivate comparing several numerical formulations: shifted SVD, shift-free SVD, banded primal factorization, dual banded systems, and suitably scaled least-squares or QR formulations. The theoretical asymptotic cost alone cannot certify their accuracy.

\paragraph{Three distinct notions of failure.}
A factorization failure means an algorithm could not produce the required numerical factors. A \emph{post-solve} failure means factors and a candidate fit existed but its residual, objective or KKT check did not satisfy the prescribed tolerance. A \emph{parameter-mapping} failure means the inversion of $S_d(\lambda)=s$ did not attain its tolerance; it is not the same as a failed trend solve. Memory exclusion, a missing optional dependency and an undefined statistical criterion are further distinct statuses. Comparing only successful runtime rows while silently dropping failures can produce misleading algorithm rankings.

For a computed finite-penalty trend $\widetilde\tau$, the first-order residual is
\begin{equation}\label{eq:kkt}
g_\lambda=y-\widetilde\tau
-\lambda D_d^\top(D_d\widetilde\tau-\mu\mathbf1_r).
\end{equation}
Its appropriately scaled norm measures a backward or stationarity residual. A small residual is necessary for trusting a fit, but can coexist with substantial forward error when $\kappa_2(A_\lambda)$ is large. Independent higher-precision reference solutions and objective comparisons are therefore useful for hard cases. The originating numerical-reliability manuscript includes archived finite-precision comparisons for high difference orders; they are preserved as supporting material rather than made the empirical center of this theory article.

\paragraph{Stable evaluation of spectral weights.}
For interior penalties, form $w_j$ from $\log\lambda+2\log\sigma_j$ with overflow-safe logistic arithmetic rather than directly forming $\lambda\sigma_j^2$ when it may overflow. The exact $s=0,1$ trends use Corollary~\ref{cor:endfits}. Likewise $S\mapsto\lambda$ should be inverted with bracketed monotone root finding, explicit convergence checks and suitable endpoint guards; a continuous inverse exists mathematically but its floating-point realization is not automatic.

\paragraph{What Cholesky failure does and does not establish.}
Since $A_\lambda$ is SPD, failure of an ordinary float64 primal Cholesky routine in high-order and high-smoothness configurations is not a counterexample to positive definiteness. It is evidence about that finite-precision formulation, scaling and implementation. The archived high-precision companion analysis reported primal banded-Cholesky failures in difficult configurations, while shift-free SVD had substantially lower measured forward error there. Those outcomes do not imply that every structured Cholesky algorithm or precision level fails, nor that SVD is globally safer for every order and penalty. Refined and dual structured methods remain admissible candidates after their own verification.

We define the valid method set for a workload by
\begin{equation}\label{eq:valid}
\mathfrak A_{\rm valid}
=\{\mathcal A:\operatorname{valid}(\mathcal A)=1\}.
\end{equation}
Here validity requires passing the memory, precision,
convergence, criterion and endpoint checks.
A nominally inexpensive method outside this set cannot be declared the computational winner.

\section{A workload-aware algorithm-selection policy}\label{sec:policy}
The theoretical problem is to choose a complete pipeline, not a favored numerical library in isolation. For a specified tolerance $\varepsilon$, memory budget $B$, difference order $d$, number of candidates $M$, observed-series count $R$ and criterion set $\mathcal C$, we recommend
\begin{equation}\label{eq:frontier}
\mathcal A^\star\in
\arg\min_{\mathcal A\in\mathfrak A_{\rm valid}}
T_\mathcal A(n,d,M,R,\mathcal C;\varepsilon,B),
\end{equation}
where the constraint \eqref{eq:valid} is evaluated before comparing cost. Neither Theorem~\ref{thm:svdwork} nor Theorem~\ref{thm:bandwork} supplies measured wall-clock constants, so \eqref{eq:frontier} is a conditional decision principle, not an empirically calibrated universal dispatcher.

\begin{table*}[t]
\caption{Algorithm choices suggested by the derived cost components. Each preference is conditional on accurate outputs, an admissible $S$ inversion and the complete criterion being available.}
\label{tab:policy}\centering\small
\setlength{\tabcolsep}{5pt}
\begin{tabularx}{\textwidth}{@{}p{.24\textwidth}Xp{.26\textwidth}@{}}
\toprule
Workload & Initial candidate and rationale & Important qualification\\
\midrule
One curve, known $\lambda$, small fixed $d$ & Banded Cholesky: $O(nd^2)$ factor, $O(nd)$ solve & Large $\lambda$ can amplify finite-precision error.\\
One curve, prescribed $S$ & Banded fit with a certified $S\mapsto\lambda$ mapper & Do not omit mapping cost. For $d=1$ the spectrum is analytic.\\
Many $S$ values, scalar GCV or likelihood, moderate $n$ & Retained SVD: costly setup, $O(r)$ candidate scores & Dense basis must fit memory; finite constants can reverse comparisons.\\
Many $S$ values, exact ordinary CV, small $d$ and large $n$ & Banded Cholesky plus selected inverse & Spectral CV also has $O(nr)$ leverage/residual reconstruction.\\
Many series sharing $(n,d,\mathcal S)$ & Map $S$ once; reuse spectrum or factors and leverage diagnostics & Each series still needs its own projections or solves.\\
Near $S=1$, large $d$, difficult conditioning & Validated shift-free spectral or independently verified structured alternative & Prefer analytic endpoint at exact $S=1$; never assume a particular algorithm always passes.\\
\bottomrule
\end{tabularx}
\end{table*}

\subsection{Complete candidate-grid procedure}
The following abstract algorithm specifies a reproducible computation rather than claiming that an optimal dispatcher is already implemented in the library.
\begin{enumerate}[leftmargin=*]
\item Fix $(n,d)$, the reference-mean policy, the criterion set $\mathcal C$, candidate set $\mathcal S$, numerical tolerances and memory limit. Declare whether any mean estimate is conditional or foldwise-refitted.
\item Determine the exact admissible endpoint conventions. Map each distinct interior $s_j$ to a penalty once, checking its inverse-map residual. Do not independently space a comparison grid in $\log\lambda$ and call it the same statistical workload.
\item Evaluate complete pipeline feasibility. A spectral method must budget its full basis; a structured method must budget all selected-inverse and likelihood diagnostics, not merely the trend fit.
\item Reuse calculations at the correct level: the $S$ mapping is shared across all series of the same length and order; a banded factor or leverage vector is shared across series for the same penalty; spectral projections depend on the observed series and fixed mean.
\item At each candidate, evaluate all requested scores and numerical checks. Record invalid or undefined scores and failed fits rather than assigning them artificial best values.
\item Return a minimizer \emph{over the explicitly evaluated finite grid} for each requested criterion. The result is not automatically a certified global minimizer of the unrestricted continuous criterion.
\end{enumerate}

\subsection{What can be concluded without empirical timing}
The theory proves that structured factorization has favorable asymptotic storage and single-fit arithmetic for fixed small $d$. It also proves that once an SVD has been paid for, many scalar smoothing scores are cheap and share the same spectral basis and smoothness map. It does \emph{not} establish that SVD always beats banded Cholesky on CV grids, that a positive crossover exists in every regime, or that a faster measured fit implies faster complete criterion selection. Those are contingent questions involving operation constants and numerical admissibility.

Finally, changing $d$ is a model-design decision, not merely a computational tuning parameter. Different difference orders impose different nullspaces and different smoothing trajectories. A preference inferred from computational work must not be misreported as a guarantee of better latent-trend recovery or forecasting accuracy.

\section{Reproducible software interface}\label{sec:software}
The \texttt{trendestimation} Python package provides an explicit user-facing interface to the smoothing problem, exposing the difference order, prescribed reference mean and normalized smoothness coordinate. The library's present accuracy-oriented default uses a retained SVD. The analytical computational policy in Section~\ref{sec:policy} should not be mistaken for an already deployed universal automatic backend selector.

A minimal example with a prescribed nonzero third-difference reference is:
\begin{verbatim}
from trendestimation import TrendSmoother
from trendestimation import select_lambda

solver = TrendSmoother(len(y), order=3)
fit = solver.fit_smoothness(
    y, smoothness=0.8, mu=0.15
)
trend = fit.trend
choice = select_lambda(
    solver, y, criterion="gcv", mu=0.15
)
selected = solver.fit(
    y, lambda_=choice.lambda_, mu=0.15
)
\end{verbatim}
The value $\mu=0.15$ is merely a demonstration of API syntax, \emph{not} a scientifically justified physical reference value. In a real study its units, source and inferential role must be defined. An alternative \texttt{mu="data"} computes an observed-difference plug-in from the supplied sample; it must not be described as fixed, genuinely foldwise-refitted CV, or joint unknown-mean likelihood.

The interface distinguishes (i) specifying a penalty, (ii) specifying normalized smoothness and internally inverting it, and (iii) selecting a candidate by a named criterion. The separate functions make the mathematical transformations inspectable and permit independent numerical QA. The library is a research implementation of established PLS smoothing methods; this paper makes no claim that no related smoothing software existed previously.

\paragraph{Reproducibility of the theory.}
The companion repository retains the source of both predecessor manuscripts, proof derivations, deterministic matrix-identity checks, and the archived numerical and computational studies. The present manuscript neither overwrites nor reruns the expensive experiments. Its independent small deterministic verification reproduces the SVD fit identity, the smoothness map, a banded Cholesky solve and selected inverse, and the determinant identity under nonzero reference means. All mathematical claims and any subsequent measured results must remain associated with exact source revisions, numerical tolerances and software dependencies.

\section{Discussion and conclusions}\label{sec:conclusion}
A bounded trace-derived smoothness coordinate turns the classical PLS penalty into an interpretable input that can be prescribed across difference orders. The normalized interval $[0,1]$ has exact interpolation and constrained-polynomial endpoints, including nonzero constant $d$th-difference reference levels. The order-$d$ singular-vector representation supplies the penalty spectrum and a shift-free reconstruction that avoids explicit large polynomial reference vectors. These facts explain why diagonalization can serve as a shared computational engine for both smoothing and smoothness-parameter selection.

The algorithmic choice is more subtle than selecting the fastest linear solver. A full SVD costs more to prepare and retain than a banded factorization, but enables inexpensive evaluation of many scalar criteria after preprojection. Banded Cholesky, when coupled with selected inversion, supplies complete conditional CV, GCV, information and likelihood quantities with structured storage and an $O(nd^2)$ per-candidate bound. The initial cost of obtaining $\lambda$ from an $S$ request and the number of observed series materially alter that comparison. Exact ordinary CV can require more spectral work than scalar GCV; no universal SVD--Cholesky winner follows from the asymptotic bounds.

The same theory also identifies a risk frontier. Large penalties and higher difference orders can create ill-conditioned normal equations, even though the matrices remain SPD mathematically. A factorization that is computationally inexpensive but fails a numerical-validity requirement does not supply a usable statistical estimator. Numerical failures must be classified by their mechanism rather than interpreted as proof that Cholesky or spectral methods are universally unsuitable.

Our principal scope restriction is a prescribed numeric $\mu$. The general nonzero-reference case is scientifically relevant and computationally accessible without changing the spectrum; jointly estimating $\mu$ and $\lambda$ is an extension that changes degrees of freedom and criterion definitions and deserves a separate inferential treatment. Similarly, optimizing the degree $d$, using ensembles, or assessing future forecast horizons requires validation objectives distinct from the retrospective trend-smoothing criteria derived here.

Whittaker--Henderson smoothing, its spectral theory, numerical factorization methods, and classical GCV were established before this work. The contribution of the present article is the integrated, unit-interval, workload-dependent mathematical and software framework and its explicit numerical validity conditions. Empirical timings and physical forecasting performance are investigated in separate companion manuscripts, rather than being inferred from theoretical operation counts.

\section*{Code and research-material availability}
The Python implementation, independently preserved Paper A and Paper B sources, theoretical verification programs, and archived supporting evidence are available at \url{https://github.com/heri-espino/trendestimation}. The canonical source is available in
\path{paper/csda-unified-theory/}; its revision
should be frozen and cited before submission. No new benchmark measurements are claimed in this article.

\section*{Declarations}
Author ordering, affiliations, corresponding-author metadata, funding, conflicts of interest and individual contributions require explicit coauthor verification before journal submission. The present draft does not assert unverified declarations.

\bibliographystyle{cas-model2-names}
% References are reused without duplicating/diverging the existing BibTeX source.
\bibliography{../csda-computational-costs/references}
\end{document}
